Hệ thống dùng 8 bảng DynamoDB. Trước khi tạo, cần hiểu một điểm khác
biệt căn bản so với cơ sở dữ liệu quan hệ.

\vspace{0.5em}\noindent\textbf{Thiết kế theo mẫu truy
vấn}\par

Với cơ sở dữ liệu quan hệ, bạn chuẩn hoá dữ liệu trước rồi viết câu truy
vấn sau. Với DynamoDB thì ngược lại: bạn xác định trước sẽ truy vấn thế
nào, rồi thiết kế bảng phục vụ đúng những truy vấn đó.

Cụ thể ở đây, sáu bảng đầu chỉ cần tra theo một khoá duy nhất. Nhưng hai
bảng cuối thì khác:

\begin{itemize}
\tightlist
\item
  \textbf{Reviews}: cần lấy \emph{tất cả đánh giá của một sản phẩm}
\item
  \textbf{Orders}: cần lấy \emph{tất cả đơn hàng của một người dùng}
\end{itemize}

Nếu chỉ dùng khoá đơn, muốn lấy các bản ghi đó phải chạy \textbf{Scan}
đọc toàn bộ bảng rồi lọc. Tốn thời gian và tốn tiền vì DynamoDB tính phí
theo lượng dữ liệu đọc.

Giải pháp là dùng \textbf{khoá tổ hợp} gồm khoá phân vùng và khoá sắp
xếp. Khi đó chỉ cần \textbf{Query} đọc đúng phân vùng cần thiết.

{\def\LTcaptype{none} % do not increment counter
\begingroup
\small
\setlength{\tabcolsep}{3pt}
\renewcommand{\arraystretch}{1.15}
\begin{longtable}[]{@{}llll@{}}
\toprule\noalign{}
Bảng & Khoá phân vùng & Khoá sắp xếp & Nội dung \\
\midrule\noalign{}
\endhead
\bottomrule\noalign{}
\endlastfoot
Products & \texttt{id} & --- & Thông tin sản phẩm \\
Categories & \texttt{id} & --- & Danh mục \\
Users & \texttt{id} & --- & Tài khoản, mật khẩu đã băm \\
Sessions & \texttt{token} & --- & Phiên đăng nhập \\
Carts & \texttt{userId} & --- & Giỏ hàng \\
Vouchers & \texttt{code} & --- & Mã giảm giá \\
Reviews & \texttt{productId} & \texttt{id} & Đánh giá sản phẩm \\
Orders & \texttt{userId} & \texttt{id} & Đơn hàng \\
\end{longtable}
\endgroup
}

\textbf{Lưu ý cho Production:} DynamoDB rất phù hợp cho các luồng giao
dịch tốc độ cao (như giỏ hàng, đơn hàng). Tuy nhiên, đối với các truy
vấn tìm kiếm và lọc sản phẩm phức tạp (như tìm theo khoảng giá, màu sắc,
từ khóa tự do), bạn nên cân nhắc đồng bộ dữ liệu sang \textbf{Amazon
OpenSearch} để xử lý thay vì phụ thuộc hoàn toàn vào DynamoDB.

\vspace{0.5em}\noindent\textbf{Bước 1. Tạo bảng bằng
console}\par

Làm thử một bảng trên console để hiểu giao diện:

\begin{enumerate}
\def\labelenumi{\arabic{enumi}.}
\tightlist
\item
  Mở
  \href{https://ap-southeast-1.console.aws.amazon.com/dynamodbv2/home?region=ap-southeast-1}{DynamoDB
  console}
\item
  Chọn \textbf{Create table}
\item
  Điền:

  \begin{itemize}
  \tightlist
  \item
    \textbf{Table name}: \texttt{Products}
  \item
    \textbf{Partition key}: \texttt{id}, kiểu \textbf{String}
  \item
    Bỏ trống Sort key
  \item
    \textbf{Table settings}: chọn \textbf{Customize settings}
  \item
    \textbf{Table class}: DynamoDB Standard
  \item
    \textbf{Capacity mode}: \textbf{On-demand}
  \end{itemize}
\item
  Chọn \textbf{Create table}
\end{enumerate}

\begin{figure}
\centering
\pandocbounded{\includegraphics[keepaspectratio,alt={Tạo bảng}]{5-Workshop/5.4/create-table.png}}
\caption{Tạo bảng}
\end{figure}

\textbf{On-demand} nghĩa là bạn trả tiền theo số lần đọc ghi thực tế,
không phải theo dung lượng đặt trước. Với dự án chưa biết lưu lượng thì
đây là lựa chọn an toàn, không sợ vừa thừa tiền vừa nghẽn.

\vspace{0.5em}\noindent\textbf{Bước 2. Tạo bảy bảng còn lại bằng
CLI}\par

Tạo bảy bảng bằng tay khá mất thời gian. Dùng lệnh sau cho nhanh:

\begin{Shaded}
\begin{Highlighting}[]
\ControlFlowTok{for}\NormalTok{ T }\KeywordTok{in}\NormalTok{ Categories Users Sessions Carts Vouchers}\KeywordTok{;} \ControlFlowTok{do}
  \ControlFlowTok{case} \VariableTok{$T} \KeywordTok{in}
    \SpecialStringTok{Sessions}\KeywordTok{)} \VariableTok{K}\OperatorTok{=}\NormalTok{token }\ControlFlowTok{;;}
    \SpecialStringTok{Carts}\KeywordTok{)}    \VariableTok{K}\OperatorTok{=}\NormalTok{userId }\ControlFlowTok{;;}
    \SpecialStringTok{Vouchers}\KeywordTok{)} \VariableTok{K}\OperatorTok{=}\NormalTok{code }\ControlFlowTok{;;}
    \PreprocessorTok{*}\KeywordTok{)}        \VariableTok{K}\OperatorTok{=}\NormalTok{id }\ControlFlowTok{;;}
  \ControlFlowTok{esac}
  \ExtensionTok{aws}\NormalTok{ dynamodb create{-}table }\AttributeTok{{-}{-}table{-}name} \VariableTok{$T} \DataTypeTok{\textbackslash{}}
    \AttributeTok{{-}{-}attribute{-}definitions}\NormalTok{ AttributeName=}\VariableTok{$K}\NormalTok{,AttributeType=S }\DataTypeTok{\textbackslash{}}
    \AttributeTok{{-}{-}key{-}schema}\NormalTok{ AttributeName=}\VariableTok{$K}\NormalTok{,KeyType=HASH }\DataTypeTok{\textbackslash{}}
    \AttributeTok{{-}{-}billing{-}mode}\NormalTok{ PAY\_PER\_REQUEST}
\ControlFlowTok{done}
\end{Highlighting}
\end{Shaded}

Hai bảng có khoá tổ hợp tạo riêng:

\begin{Shaded}
\begin{Highlighting}[]
\ExtensionTok{aws}\NormalTok{ dynamodb create{-}table }\AttributeTok{{-}{-}table{-}name}\NormalTok{ Reviews }\DataTypeTok{\textbackslash{}}
  \AttributeTok{{-}{-}attribute{-}definitions}\NormalTok{ AttributeName=productId,AttributeType=S AttributeName=id,AttributeType=S }\DataTypeTok{\textbackslash{}}
  \AttributeTok{{-}{-}key{-}schema}\NormalTok{ AttributeName=productId,KeyType=HASH AttributeName=id,KeyType=RANGE }\DataTypeTok{\textbackslash{}}
  \AttributeTok{{-}{-}billing{-}mode}\NormalTok{ PAY\_PER\_REQUEST}

\ExtensionTok{aws}\NormalTok{ dynamodb create{-}table }\AttributeTok{{-}{-}table{-}name}\NormalTok{ Orders }\DataTypeTok{\textbackslash{}}
  \AttributeTok{{-}{-}attribute{-}definitions}\NormalTok{ AttributeName=userId,AttributeType=S AttributeName=id,AttributeType=S }\DataTypeTok{\textbackslash{}}
  \AttributeTok{{-}{-}key{-}schema}\NormalTok{ AttributeName=userId,KeyType=HASH AttributeName=id,KeyType=RANGE }\DataTypeTok{\textbackslash{}}
  \AttributeTok{{-}{-}billing{-}mode}\NormalTok{ PAY\_PER\_REQUEST}
\end{Highlighting}
\end{Shaded}

\vspace{0.5em}\noindent\textbf{Bước 3. Kiểm tra}\par

\begin{Shaded}
\begin{Highlighting}[]
\ExtensionTok{aws}\NormalTok{ dynamodb list{-}tables}
\end{Highlighting}
\end{Shaded}

Kết quả phải liệt kê đủ 8 bảng.

\vspace{0.5em}\noindent\textbf{Bước 4. Nạp dữ liệu
mẫu}\par

Nạp danh mục và sản phẩm mẫu để hệ thống có thứ hiển thị:

\begin{Shaded}
\begin{Highlighting}[]
\BuiltInTok{cd}\NormalTok{ backend}
\ExtensionTok{node}\NormalTok{ scripts/seed{-}data.mjs}
\end{Highlighting}
\end{Shaded}

Kiểm tra lại số bản ghi trong bảng Products:

\begin{Shaded}
\begin{Highlighting}[]
\ExtensionTok{aws}\NormalTok{ dynamodb scan }\AttributeTok{{-}{-}table{-}name}\NormalTok{ Products }\AttributeTok{{-}{-}select}\NormalTok{ COUNT}
\end{Highlighting}
\end{Shaded}
